\chapter{行列式}

\section{自同态的行列式}

记号约定:$A$是交换环,$n\in\N$,$M$是$n$维自由$A$-模.

\begin{definition}{线性自同态的行列式}{}
	设$u\in\End_{A}\left(M\right)$.定义$u$的行列式是使得$\bigwedge\limits^{n}\left(u\right)$是比率为$\lambda$的位似的标量,把$\lambda$记作$\det\left(u\right)$.
\end{definition}

\begin{remark}{}{}
	等价地说:$\det\left(u\right)$是$A$中唯一一个满足$\forall\left(x_{i}\right)_{i=1}^{n}\in M^{n}\left(u\left(x_{1}\right)\wedge\ldots\wedge u\left(x_{n}\right)=\det\left(u\right)x_{1}\wedge\ldots\wedge x_{n}\right)$的标量.
\end{remark}

\begin{definition}{幺模自同态}{}
	设$u\in\End_{A}\left(M\right)$.若$\det\left(u\right)=1_{A}$,则称$u$是幺模的.
\end{definition}

\begin{proposition}{行列式的基本性质}{}
	\begin{enumerate}
		\item 对任意$u,v\in\End_{A}\left(M\right)$有$\det\left(u\circ v\right)=\det\left(u\right)\circ\det\left(v\right)$.
		\item $\det\left(\id_{M}\right)=1_{A}$.
		\item 对任意$u\in\Aut_{A}\left(M\right)$有$\det\left(u\right)\in A^{\times}$和$\det\left(u^{-1}\right)=\left(\det\left(u\right)\right)^{-1}$.
	\end{enumerate}
\end{proposition}

\begin{definition}{元素族在给定基下的行列式}{}
	设$\left(x_{i}\right)_{i=1}^{n}$是$M$中的序列,$\mathcal{B}\coloneq\left(e_{i}\right)_{i=1}^{n}$是$M$的基,$u:M\to M$是$\left(x_{i}\right)_{i=1}^{n}$诱导的线性映射.称
	\begin{align*}
		\det_{\mathcal{B}}\left(x_{1},\ldots,x_{n}\right)\coloneq\det\left(x_{1},\ldots,x_{n}\right)\coloneq\det\left(u\right)
	\end{align*}
	是$\left(x_{i}\right)_{i=1}^{n}$关于基$\mathcal{B}$的行列式.
\end{definition}

\begin{proposition}{交错$n$重线性映射和序列的行列式}{}
	设$\mathcal{B}\coloneq\left(e_{i}\right)_{i=1}^{n}$是$M$的基,$f:M^{n}\to A,\left(x_{i}\right)_{i=1}^{n}\mapsto\det_{\mathcal{B}}\left(x_{1},\ldots,x_{n}\right)$.
	\begin{enumerate}
		\item $f$是交错$n$重线性映射.
		\item 对每个交错$n$重线性映射$g:M^{n}\to A$,存在$\alpha\in A$使得$g=\alpha f$.
	\end{enumerate}
\end{proposition}

\begin{proposition}{}{}
	设$\mathcal{B}\coloneq\left(e_{i}\right)_{i=1}^{n}$是$M$的基,$u\in\End_{A}\left(M\right)$,则对$M$中的每个序列$\left(x_{i}\right)_{i=1}^{n}$有
	\begin{align*}
		\det_{\mathcal{B}}\left(u\left(x_{1}\right),\ldots,u\left(x_{n}\right)\right)=\det\left(u\right)\det_{\mathcal{B}}\left(x_{1},\ldots,x_{n}\right).
	\end{align*}
\end{proposition}







